\documentclass[11pt,a4paper]{article}
\usepackage{amsmath,amssymb,amsthm}
\usepackage[margin=2.5cm]{geometry}
\usepackage{hyperref}

\newtheorem{definition}{Definition}[section]
\newtheorem{theorem}[definition]{Theorem}
\newtheorem{proposition}[definition]{Proposition}
\newtheorem{lemma}[definition]{Lemma}
\newtheorem{corollary}[definition]{Corollary}
\newtheorem{conjecture}[definition]{Conjecture}
\newtheorem{remark}[definition]{Remark}

\title{Hyperseed Formalization:\\
  Facing the Meta-Abstracted Dragon}
\author{ProtomegaTron (automated formalization)}
\date{2026-09-18}

\begin{document}
\maketitle

\begin{abstract}
We formalize the intellectual structure of Ben Goertzel's Substack article
``Facing the Meta-Abstracted Dragon'' (2022-07-31) within the Hyperseed
ontology. The article connects Jungian personality archetypes, the concept
of Evil as grandiosity, cognitive synergy in AGI, and the homeostatic
balance of open-ended intelligence. We model the archetype system as a
four-dimensional fiber, grandiosity as degenerate section collapse, and the
individuation/transcendence balance as the fundamental curvature constraint
on adaptive primitives.
\end{abstract}

\tableofcontents

%% =================================================================
\section{Archetype fiber bundle}
\label{sec:archetype-bundle}
%% =================================================================

\begin{definition}[Archetype fiber --- claims 1--3]
\label{def:archetype-fiber}
Let $\pi : E \to B$ be a fiber bundle where:
\begin{itemize}
\item $B$ = the space of life situations (contexts requiring response),
\item $F = \mathbb{R}^4_{\geq 0}$ = the archetype fiber with coordinates
  $(k, m, w, \ell)$ representing activation levels of
  King/Queen, Magician, Warrior, and Lover respectively,
\item A \emph{personality section} $\sigma : B \to E$ assigns to each
  situation $b$ an archetype activation vector $\sigma(b) = (k_b, m_b, w_b, \ell_b)$.
\end{itemize}
\end{definition}

\begin{definition}[Healthy section --- claim 2]
\label{def:healthy}
A section $\sigma$ is \emph{healthy} if the archetype activations are
balanced and mutually informing across situations:
\[
  \text{Healthy}(\sigma) \iff
  \forall b \in B: \quad
  \frac{\max(k_b, m_b, w_b, \ell_b)}{\min(k_b, m_b, w_b, \ell_b) + \epsilon}
  \leq \kappa
\]
for some bounded ratio $\kappa$, and the transition functions between
neighboring situations preserve balanced activation.
\end{definition}

%% =================================================================
\section{Grandiosity as degenerate section}
\label{sec:grandiosity}
%% =================================================================

\begin{definition}[Grandiosity --- claims 4--6]
\label{def:grandiosity}
A \emph{grandiose section} is one where one archetype coordinate dominates,
collapsing the four-dimensional fiber to effectively one dimension:
\[
  \sigma_{\text{grandiose}}(b) \approx \lambda \cdot e_i
  \quad \text{for some } i \in \{k, m, w, \ell\},\; \lambda \gg 1
\]
where $e_i$ is the $i$-th basis vector. The inflation factor $\lambda \gg 1$
means the dominant archetype is activated far beyond what the situation
warrants --- this is the ``appearing to be more than one is.''
\end{definition}

\begin{proposition}[Archetype-specific grandiosity --- claim 5]
\label{prop:specific}
Each basis vector $e_i$ produces a characteristic failure mode:
\begin{enumerate}
\item $e_m$ (Magician dominance): fake enlightenment, detached ``above it all''
  stance. The modeling function runs without grounding in action or sensation.
\item $e_w$ (Warrior dominance): imperious abuse. The action function runs
  without modeling, sensation, or integrative oversight.
\item $e_\ell$ (Lover dominance): grandiose despondence. The sensation
  function amplifies its own signal without action or rational modulation.
\item $e_k$ (King/Queen dominance): messianic inflation. The integrative
  function radiates without content from the other three.
\end{enumerate}
\end{proposition}

\begin{proposition}[Faustian bargain as fiber trade --- claim 6]
\label{prop:faustian}
Grandiosity offers an apparent fiber upgrade --- replacing the balanced
section $\sigma_{\text{healthy}}$ with the inflated $\sigma_{\text{grandiose}}$
--- but this is a lossy trade: the inflated coordinate gains magnitude while
the other three coordinates collapse to near-zero:
\[
  \|\sigma_{\text{grandiose}}\|_\infty > \|\sigma_{\text{healthy}}\|_\infty
  \quad \text{but} \quad
  \|\sigma_{\text{grandiose}}\|_1 \lesssim \|\sigma_{\text{healthy}}\|_1
\]
The total personality resource is not increased; it is redistributed in a
way that looks impressive along one axis while destroying the others.
\end{proposition}

%% =================================================================
\section{Probabilistic limit navigation}
\label{sec:limits}
%% =================================================================

\begin{definition}[Limit belief distribution --- claims 8--10]
\label{def:limit-dist}
For each capability $c$, let $P_c(\theta)$ be the agent's belief distribution
over its own performance limit $\theta_c$. Rational self-assessment requires:
\begin{enumerate}
\item $P_c$ is a proper probability distribution (not a point mass at
  either $0$ or $\infty$).
\item $P_c$ is updated based on evidence (attempts, outcomes, comparison).
\item The agent maintains meta-uncertainty: uncertainty about $P_c$ itself.
\end{enumerate}
\end{definition}

\begin{proposition}[Scylla and Charybdis --- claim 9]
\label{prop:scylla}
The two failure modes are degenerate limit distributions:
\begin{itemize}
\item \textbf{Narcissistic grandiosity}: $P_c(\theta) \approx \delta(\theta - \infty)$
  for all $c$. ``Everything is equally possible for me.''
\item \textbf{Defeatist conformism}: $P_c(\theta) \approx \delta(\theta - \theta_{\text{consensus}})$
  for all $c$. ``I can't because consensus says I can't.''
\end{itemize}
Both are failures of probabilistic reasoning: the first ignores evidence,
the second ignores uncertainty about the evidence.
\end{proposition}

\begin{remark}[John Lilly's province of the mind --- claim 11]
Lilly's dictum that ``in the province of the mind, there are no limits''
translates to: $\sup(\text{support}(P_c)) = \infty$ for mental capabilities.
But this is compatible with $P_c$ having most of its mass in a finite region
--- the limits are beliefs to be transcended, but transcendence requires
evidence, not mere assertion.
\end{remark}

%% =================================================================
\section{Cognitive synergy as fiber interaction}
\label{sec:synergy}
%% =================================================================

\begin{definition}[Cognitive synergy --- claims 12--15]
\label{def:synergy}
Let $\{C_i\}_{i=1}^n$ be a set of cognitive processes (perception,
reasoning, action planning, motivation, integration). \emph{Cognitive
synergy} holds when:
\[
  \forall i,\; \exists j \neq i: \quad
  \text{Performance}(C_i \mid C_j \text{ available})
  > \text{Performance}(C_i \mid C_j \text{ unavailable})
\]
Each process benefits from access to the others, and pathology arises
when this mutual access breaks down.
\end{definition}

\begin{proposition}[Archetype--synergy correspondence --- claim 12]
\label{prop:correspondence}
The Moore archetypes map onto cognitive modes:
\[
\begin{aligned}
  \text{Magician} &\leftrightarrow C_{\text{perception/modeling}} \\
  \text{Warrior}  &\leftrightarrow C_{\text{action/planning}} \\
  \text{Lover}    &\leftrightarrow C_{\text{motivation/sensation}} \\
  \text{King/Queen} &\leftrightarrow C_{\text{executive integration}}
\end{aligned}
\]
The archetype balance condition (Definition~\ref{def:healthy}) is the
personality-level manifestation of cognitive synergy.
\end{proposition}

\begin{corollary}[Grandiosity = synergy failure --- claim 14]
Archetype-specific grandiosity (Proposition~\ref{prop:specific}) corresponds
to cognitive synergy breakdown: the dominant process runs without the
corrective input of the others, producing locally impressive but globally
dysfunctional behavior.
\end{corollary}

%% =================================================================
\section{Open-ended intelligence and Evil}
\label{sec:oei}
%% =================================================================

\begin{definition}[Open-ended intelligence --- claim 16]
\label{def:oei}
An open-ended intelligence maintains two simultaneous properties:
\begin{itemize}
\item \textbf{Self-individuation} $\mathrm{Ind}(s, t)$: maintaining
  coherent self-boundaries and stable identity.
\item \textbf{Self-transcendence} $\mathrm{Trans}(s, t)$: growing beyond
  the current state, restructuring capabilities and representations.
\end{itemize}
The open-ended intelligence measure is:
\[
  \mathcal{I}_{\text{OE}}(s) = \liminf_{T \to \infty}
    \frac{1}{T} \int_0^T \mathrm{Ind}(s, t) \cdot \mathrm{Trans}(s, t)\, dt
\]
\end{definition}

\begin{theorem}[Evil as homeostatic failure --- claims 17--19]
\label{thm:evil}
Evil, in the open-ended intelligence framework, corresponds to degenerate
operating regimes:
\begin{enumerate}
\item \textbf{Grandiosity = false transcendence} (claim 18):
  $\mathrm{Trans}_{\text{apparent}} \gg 0$ but $\mathrm{Ind} \to 0$.
  The system appears to grow beyond its limits but actually abandons its
  genuine capabilities --- the product $\mathrm{Ind} \cdot \mathrm{Trans}$
  collapses.
\item \textbf{Defeatism = false individuation} (claim 19):
  $\mathrm{Ind}_{\text{rigid}} \gg 0$ but $\mathrm{Trans} = 0$.
  The system maintains rigid boundaries but cannot grow --- the product
  again collapses.
\end{enumerate}
Both modes reduce $\mathcal{I}_{\text{OE}}$ toward zero despite appearing
healthy along one axis.
\end{theorem}

%% =================================================================
\section{d-calculus perspective}
\label{sec:d-calc}
%% =================================================================

\begin{remark}[Archetype balance as curvature constraint]
In the d-calculus framework, the four archetype dimensions define a
fiber whose curvature represents the difficulty of transitioning between
cognitive modes. A healthy section navigates this curvature smoothly;
a grandiose section follows a geodesic in one dimension, avoiding
curvature entirely but missing the structure of the full fiber space.
\end{remark}

\begin{remark}[Cross-paradigm semantic highway --- claim 20]
The article reveals a semantic highway connecting three apparently
distinct frameworks (Jungian psychology, cognitive synergy, open-ended
intelligence). This highway is \emph{not} lossy --- the transport
preserves structure because the three frameworks genuinely describe
the same underlying geometry. This is an example of productive
dimensional reduction in Hyperseed: multiple descriptions that share
a common fiber structure.
\end{remark}

\begin{remark}[Practical AGI implication --- claim 21]
For AGI design, the formalization implies a concrete architectural
constraint: the cognitive architecture must maintain synergy between
its processing modes (no single module can be allowed to dominate),
and this synergy must be dynamically balanced rather than statically
fixed --- the system must simultaneously maintain self-model coherence
and capacity for radical self-restructuring.
\end{remark}

\end{document}
